\documentclass[12pt, a4paper]{report}

% Paquetes básicos
\usepackage[utf8]{inputenc} 
\usepackage[spanish]{babel} 
\usepackage{graphicx} 
\usepackage{hyperref} 
\usepackage{geometry}
\geometry{a4paper, margin=2.5cm}

% Información del documento
\title{\textbf{Memoria Técnica: SmartPresence}}
\author{Angel Sanz Barragan \and Mark Ruiz Correa}
\date{Sprint 1 - 21/10/2026}

\begin{document}

\maketitle
\tableofcontents
\newpage

\chapter{Introducción, problema y objetivos}

\section{Nombre del proyecto y frase de venta}
El proyecto se denomina \textbf{SmartPresence}. \\
Nuestra frase de venta que define el producto es: \textit{"La digitalización definitiva del control horario del profesorado."}

\section{Problema o necesidad detectada y usuario/cliente objetivo}
\subsection{Problema detectado}
En la actualidad, muchos centros educativos siguen dependiendo de sistemas manuales o plataformas informáticas lentas para registrar la asistencia del profesorado. Esto genera un control ineficiente, pérdida de tiempo para los docentes y dificultad para que el equipo directivo sepa en tiempo real qué profesores se encuentran dentro del recinto escolar ante cualquier incidencia o sustitución.

\subsection{Usuario y cliente objetivo}
El sistema está diseñado para dos perfiles principales:
\begin{itemize}
    \item \textbf{Cliente (Equipo Directivo / Administración):} Quienes necesitan visualizar el estado de presencia en tiempo real de la plantilla para gestionar sustituciones y emergencias.
    \item \textbf{Usuario Final (Profesorado de ciclos formativos DAM/ASIX):} Quienes interactúan con el hardware para fichar de manera rápida y sin fricciones mediante su tarjeta RFID.
\end{itemize}

\section{Objetivos generales del proyecto}
Para resolver el problema planteado, el proyecto SmartPresence persigue los siguientes objetivos generales:
\begin{itemize}
    \item \textbf{Automatización del registro:} Diseñar e implementar un dispositivo de hardware IoT basado en un microcontrolador ESP32 y un lector MFRC522 que permita el fichaje autónomo mediante tarjetas RFID.
    \item \textbf{Centralización de datos:} Desarrollar una API REST robusta que actúe como backend, soportada por una base de datos relacional desplegada en contenedores Docker, garantizando la persistencia y seguridad de la información.
    \item \textbf{Accesibilidad móvil:} Crear una aplicación móvil multiplataforma utilizando el framework Flutter, que consuma la API y muestre de forma visual y en tiempo real (mediante indicadores semafóricos verde/rojo) la presencia del profesorado en el centro.
    \item \textbf{Aplicación de metodologías ágiles:} Gestionar todo el ciclo de vida del desarrollo de software mediante Scrum, versionado de código en GitHub y documentación técnica en formato LaTeX.
\end{itemize}

\end{document}
